% TeX program = lualatex
% !TeX encoding = UTF-8

% NOTE: For refrence https://www.iist.ac.in/sites/default/files/2025-05/SC14D003%20%281%29.pdf and 
% - https://www.iist.ac.in/posters-presented-students-conferences-symposia 


\documentclass[20pt]{extarticle}

\usepackage[T1]{fontenc}
\usepackage{amssymb} % unicode-math handles this
\usepackage{fontspec}
\setmainfont{TeX Gyre Heros} % sans-serif
\usepackage{unicode-math}
\setmathfont{STIX Two Math}  % scalable math font

% \usepackage[framemethod=TikZ]{mdframed}
% \usepackage{pgfplots}
% \usepackage{utopia}
% \usepackage{changepage}
% \usepackage{moresize}
% \usepackage{mathpazo}
% \usepackage{tikz-3dplot}
% \usepackage{cancel}
% \tdplotsetmaincoords{70}{165}
% \pgfplotsset{compat=1.18}

\usepackage{amsthm}
\usepackage{amsmath}
\usepackage{fancyhdr}
\usepackage{flushend}

% Settings
\usepackage[a0paper, portrait, margin=2cm]{geometry}
\setlength{\columnsep}{0.2cm}  % Space between columns
\setlength{\parskip}{0cm}
\setlength{\parindent}{0cm}
\setlength{\topskip}{0cm}
\setlength{\partopsep}{0cm}
\linespread{1.2}\selectfont
\newlength{\footerheight}
\setlength{\footerheight}{3cm}

\usepackage{multicol}
\setlength{\multicolsep}{0.2cm}  % Space above/below multicols
\setlength{\premulticols}{0pt}
\setlength{\postmulticols}{0.2pt}

\usepackage{enumitem}
% \setlist{nosep, leftmargin=*}
\newcommand{\hollowsq}{\tikz[baseline=-0.5ex]\draw[TitleBlue, thick] (0,0) rectangle (0.3cm,0.3cm);}
\newcommand{\shadowsq}{%
    \begin{tikzpicture}[baseline=-0.5ex]
        \fill[gray!90] (0.05,-0.05) rectangle (0.3,0.2);
        \fill[black] (0,0) rectangle (0.25,0.25);
        % \draw[white, very thin] (0,0.25) -- (0,0) -- (0.25,0);
    \end{tikzpicture}%
}
\setlist{label=\shadowsq, itemsep=0.1pt, parsep=0pt, topsep=1.5pt, leftmargin=30pt}

\usepackage{tikz}
\usetikzlibrary{
    shadings, calc, patterns, angles, quotes, arrows.meta, fadings, 3d, perspective, backgrounds, intersections, bending,
    decorations.pathmorphing, decorations.pathreplacing, decorations.markings
}

\usepackage{xcolor}
% \definecolor{TitleBlue}{HTML}{003366}
\definecolor{TitleBlue}{HTML}{365794}
\definecolor{White}{HTML}{FFFFFF}
% \definecolor{BoxBg}{HTML}{E4DDC9}
\definecolor{BoxBg}{HTML}{F8F9F9}
\definecolor{Text}{HTML}{1A1A1A}
\definecolor{LightBlue}{RGB}{191, 223, 255}

\usepackage{caption}
\DeclareCaptionFont{postercaption}{\fontsize{25}{27}\selectfont}
\captionsetup{
	font=postercaption,
	labelfont={postercaption,bf,color=TitleBlue},
	labelsep=period,
	margin=20pt,
	justification=centering,
	hypcap=false
}

% \usepackage[many]{tcolorbox}
\usepackage[most]{tcolorbox}
\tcbset{
	arc=5pt,
	before skip=5pt, after skip=5pt,
	left=20pt, right=20pt, top=25pt, bottom=20pt,
	boxrule=3pt,
}

\newtcolorbox{conceptbox}[2][]{
	breakable, enhanced,
	% drop shadow,
	% drop lifted shadow,
	% drop shadow={2mm}{-2mm}{0mm}{TitleBlue!50!white}, % {x-offset}{y-offset}{spread}{color}
	drop fuzzy shadow={black!50!white},
	colback=BoxBg, colframe=TitleBlue, coltitle=White,
	fonttitle=\bfseries\fontsize{35}{38}\selectfont,
	% title=\MakeUppercase{#2},
	title=#2,
	before upper={\fontsize{27}{30}\selectfont\color{Text}}
}

\newcommand{\paraseperator}[2][TitleBlue]{%
\begin{center}
	\begin{tikzpicture}
		\node (textnode) [text=#1, anchor=center, inner sep=15pt] at (0,0)
		{\fontsize{25}{26}\selectfont\textbf{#2}};
		\draw [draw=#1, ultra thick] (-0.48\linewidth,0) -- (textnode.west);
		\draw [draw=#1, ultra thick] (textnode.east) -- (0.48\linewidth,0);
	\end{tikzpicture}
\end{center}
\vspace{-0.8cm}
}

% footer
\usepackage{eso-pic}  % For absolute positioning
\newcommand{\posterfooter}[1]{%
    % \begin{center}
    %     \fontsize{20}{22}\selectfont\textcolor{gray}{#1}
    % \end{center}
	\begin{tikzpicture}[overlay,remember picture]
	    \node[anchor=south, text=gray, font=\fontsize{20}{22}\selectfont] at ([yshift=2cm] current page.south){#1};
	\end{tikzpicture}%
}

\usepackage[
  pdfauthor={Harshit Prashant Dhanwalkar},
  pdftitle={Poster},
  pdfsubject={MTech Poster presentation},
  pdfkeywords={Monsoon Depressions, Brown Ocean Effect, Soil Moisture, Total Precipitation, ERA5},
  pdfcreator={LuaLaTeX},
  pdfproducer={LuaTeX-1.14.0},
  pdfpagelayout={SinglePage},
  pdfdisplaydoctitle=true,
]{hyperref}

\pagestyle{empty}
\makeatletter
\renewcommand{\@oddfoot}{}
\renewcommand{\@evenfoot}{}
\makeatother

\begin{document}
% Poster header
% \begin{tcolorbox}[colback=TitleBlue, colframe=TitleBlue, arc=0pt, sharp corners]
\begin{tcolorbox}[colback=TitleBlue, colframe=TitleBlue, arc=10pt, boxrule=0pt]
	\begin{minipage}[c]{0.09\textwidth}
		\hfill
	\end{minipage}
	\begin{minipage}[c]{0.8\textwidth}
		\centering
		\color{white}
		\vspace{0.2cm}
		{\fontsize{70}{90}\selectfont \textbf{Monsoon Depressions}} \\[0.4cm]
		{\fontsize{32}{38}\selectfont Harshit P. Dhanwalkar, A. Chandrasekar} \\[0.3cm]
		{\fontsize{20}{21}\selectfont Email-id: harshit.sc21b164@iist.ac.in, chandra@iist.ac.in} \\[0.3cm]

		{\fontsize{25}{30}\selectfont Department of Earth and Space Sciences (ESS)} \\[0.2cm]
		{\fontsize{25}{30}\selectfont Indian Institute of Space Science and Technology, Thiruvananthapuram, India}
	\end{minipage}
	\begin{minipage}[c]{0.09\textwidth}
		\centering
		\includegraphics[width=0.95\linewidth]{images/iist-logo-square.png}
	\end{minipage}
\end{tcolorbox}

% Poster footer
\AddToShipoutPictureBG*{%
	\posterfooter{Presented at: IIST Symposium, 26 March 2026}
}

\vspace{0.5cm}
\begin{conceptbox}{\centering Abstract}
	% FIX: not all though .... some desepates before.
	This study investigates the post-landfall sustainability of 2024 Monsoon Depressions (MDs) over the Indian subcontinent. Using ERA5 reanalysis and high-resolution tracking, we analyze the feedback between soil moisture (\texttt{SWVL1}) and precipitation (\texttt{TP}). Preliminary results suggest a \textbf{Brown Ocean Effect}, where antecedent soil moisture acts as a latent heat reservoir, offsetting typical inland decay and extending MD longevity.
\end{conceptbox}

\begin{multicols*}{2}
	\begin{conceptbox}{Introduction}
		\#TODO \\
		Monsoon Depressions are synoptic-scale weather systems that form during the Indian monsoon and contribute significantly to the total rainfall. While their formation and oceanic evolution are well-studied, their behavior after landfall remains poorly understood. This work examines the 2024 depressions to identify factors controlling their inland sustainability.
	\end{conceptbox}

	\begin{conceptbox}{Methodology}
		\#TODO
		\paraseperator{Tracking \& Detection}
		MD centers were identified using \texttt{MSLP} minima and 850 hPa vorticity.% Translational speeds were computed using Great Circle distances.
		% FIX:

		\paraseperator{Azimuthal Averaging}
		To analyze radial evolution, variables were averaged across 16 directional sectors in concentric annuli (0-400 km) around the MD center:
		\begin{itemize}
			\item \textbf{Disk 0:} Central point interpolation.
			\item \textbf{Annuli:} 25km increments with $\pm$10km tolerance.
		\end{itemize}
	\end{conceptbox}

	\begin{conceptbox}{Observations}
		\#TODO
		\paraseperator{Diurnal Variation}

		\paraseperator{Surface Heat Flux Transitions}
		% Latent and Sensible heat fluxes were de-accumulated from ERA5 ($J\,m^{-2}$ to $W\,m^{-2}$).
		% \[ Q = -\frac{\Phi_{t} - \Phi_{t-1}}{\Delta t} \]

		\paraseperator{Landfall Impact}
		Analysis of 2024 tracks shows that for Long-lived MDs ($\geq 96h$), surface fluxes do not drop as sharply as expected upon landfall, indicating a sustained energy supply from saturated soils.

		\paraseperator{Land-Surface Feedback}
		Performed Pearson and Spearman correlation analysis between \texttt{TP} and \texttt{SWVL1} specifically for land-phase points (\texttt{SWVL1} $> 0.03$).
		\begin{itemize}
			\item \textbf{Feedback Loop:} High positive correlation between land-phase precipitation and Layer-1 soil moisture.
			\item \textbf{Significance:} $p$-values indicate local surface forcing is a dominant driver for inland intensification in specific 2024 cases.
		\end{itemize}

		\paraseperator{Speed Analysis}
		Typical translational velocities were found in the range of \textbf{2-10 m/s}% (7--22 km/h).
	\end{conceptbox}

	\columnbreak
	\begin{conceptbox}{Results}
		\#TODO
		\begin{center}
			\includegraphics[width=0.9\linewidth]{../../version21-with-no-ocean-masking/output_2024/MD_track_2024.png} \\
			\captionof{figure}{Detected Monsoon Depressions for the year 2024}
			\label{fig:md_tracks}
		\end{center}

		\paraseperator{Conclusion}
		The 2024 MDs exhibited extended inland longevity primarily due to:
		\begin{itemize}
			\item High antecedent soil moisture.
			\item Ongoing precipitation maintaining the latent heat reservoir.
			\item Slower translation speeds allowing longer interaction with local moisture sources.
		\end{itemize}
	\end{conceptbox}

	\begin{conceptbox}[colback=white, fonttitle=\bfseries]{Refrences}
		\begin{itemize}
			\item \#TODO
			\item \#TODO
			\item \#TODO
			\item \#TODO
		\end{itemize}
	\end{conceptbox}
\end{multicols*}
\end{document}
